\section{Arithmetic of limits and density readers}\label{sec:arithmetic}
The arithmetic claims depend on the representation of a \emph{uniform}
family. For $X=0^{(m-1)}$, uniformly $\mathsf D_m$ data mean that an
$X$-oracle algorithm approximates $x_n$ within $2^{-p}$ uniformly in
$(n,p)$. Uniform $\mathsf S_m$ data have increasing rational
$X$-computable names, and uniform $\mathsf P_m$ data have decreasing
ones. For an individual real these agree with the one-sided cut classes,
with a separate rational-endpoint case. They do not imply that exact
strict lower cuts are uniformly decidable at the same level.

For example let $x_e=2^{-t}$ if machine $e$ first halts at step $t\geq1$,
and let $x_e=0$ otherwise. Running for $s$ steps and returning the value
found, or zero, gives error at most $2^{-s}$. This is a uniformly
computable family, but $\{e:x_e>0\}$ is the halting set. A theorem about
names cannot silently be replaced by a uniform exact-comparison theorem.

\begin{proposition}[Limit calculus]\label{prop:calculus}
For uniform names, finite outputs, and the stated existing limits, the
following are arithmetic ceilings. Every entry is sharp for abstract
data; physical attainment is asserted separately.
\begin{center}\small
\begin{tabular}{c|ccc}
Input & limit & lower limit & upper limit\\\hline
$\mathsf D_m$ & $\mathsf D_{m+1}$ & $\mathsf S_{m+1}$ & $\mathsf P_{m+1}$\\
$\mathsf S_m$ & $\mathsf P_{m+1}$ & $\mathsf S_{m+2}$ & $\mathsf P_{m+1}$\\
$\mathsf P_m$ & $\mathsf S_{m+1}$ & $\mathsf S_{m+1}$ & $\mathsf P_{m+2}$
\end{tabular}
\end{center}
Infimum and supremum send $\mathsf D_m$ data respectively to
$\mathsf P_m,\mathsf S_m$. For $\mathsf S_m$ data they give
$\mathsf P_{m+1},\mathsf S_m$; for $\mathsf P_m$ data they give
$\mathsf P_m,\mathsf S_{m+1}$. A limit with an effective error bound
preserves the input class. Computable continuous increasing maps
preserve the side; decreasing maps exchange it, on their domains.
\end{proposition}
\begin{proof}
Work relative to $X$. For a convergent computable rational sequence an
$X'$ search finds arbitrarily accurate tail-Cauchy intervals. Oracle
Cauchy names reduce asymptotic limits to rational sequences by errors
$2^{-n}$. Moreover
\[
 q<\liminf_n a_n\ \Longleftrightarrow\
 \exists r>q\ \exists N\ \forall n\geq N:\ a_n\geq r.
\]
This is a $\Sigma^X_2$ lower cut; the upper endpoint is dual. The
rational-endpoint distinction avoids incorrectly complementing strict
cuts at equality. For increasing names $a_{n,s}$, the analogous lower
cut of the liminf has an additional existential approximation witness
after the universal $n$ quantifier, hence $\Sigma^X_3$. The strict upper
cut of their limsup is $\Sigma^X_2$, giving a $\mathsf P_{m+1}$ real.
An existing limit is its own limsup. Negation proves the other row.

Dovetailing preserves suprema of left-c.e. and infima of right-c.e.
names. An infimum of left-c.e. values has a $\Sigma^X_2$ upper cut:
there is an index and a rational buffer below the tested threshold
bounding every increasing approximant. The dual gives the supremum
rule. If $|x_i-x|\leq2^{-i}$, use
$x=\sup_i(x_i-2^{-i})=\inf_i(x_i+2^{-i})$ to preserve one-sided names.
Computably continuous monotone maps preserve or reverse these names on
their domains. For a positive two-sided name, search for an interval
bounded away from zero; for a positive increasing name, search for a
positive approximant. For a decreasing name $q_s$ converging to a positive
value, each $q_s$ is positive, so apply logarithms or reciprocals directly
to these upper approximants and take monotone rational roundings. This
requires no computable positive lower bound on the limit.

Sharpness of the two-sided row is the relativized real hierarchy
\cite{zhengWeihrauch,density}. For the left-c.e. limit and limsup entries,
write a right-c.e.-in-$X'$ value as the limsup of an $X$-computable
sequence, and take its tail suprema, which are left-c.e. in $X$.
Sharpness of the left-c.e. liminf entry follows from the characterization
of lower densities of c.e. sets \cite[Theorem~5.7]{density}, relativized
to $X$; their finite prefix densities have increasing $X$-names.
Negation handles the dual entries. For an infimum target with
$\Sigma^X_2$ upper cut, write its membership as
$\exists a\forall b\,R(q,a,b)$. An increasing name starts at $q$ and
jumps to a common upper bound on the first failed witness. The infimum
over $(q,a)$ is the target. Supremum is dual. Bounded targets can be
handled in one fixed compact rational interval.
\end{proof}

\begin{lemma}[Relative density realization]\label{lem:relative-density}
For $X=0^{(k-1)}$ and $0\leq a\leq b\leq1$ in
$\mathsf S_{k+1}\times\mathsf P_{k+1}$, an $X$-computable set $A$
has lower density $a$ and upper density $b$. When $a=b$, its density
is that $\mathsf D_{k+1}$ real. A finite prescribed prefix can be added.
\end{lemma}
\begin{proof}
The relativized sequence characterizations
\cite[Theorem~5.1 and Corollaries~5.4--5.5]{density} provide a rational
$X$-computable sequence with those endpoints: for unequal endpoints,
interleave lower and upper sequences clipped on opposite sides of a
fixed rational separator. For equal endpoints use a convergent sequence.
Clip all values to $[0,1]$. On $[N_s,N_{s+1})$, where $N_0=1$ and
$N_{s+1}=(s+2)N_s$, write the mechanical word of slope $b_s$.
Its first $t$ letters contain $\lfloor t b_s\rfloor$ ones. At a block
endpoint its density differs from $b_s$ by at most
$1/(s+2)+1/N_{s+1}$. Between endpoints it differs by at most $1/J$
from a convex combination of the preceding density and $b_s$.
Thus no new extreme limit values are introduced, and both endpoints
are attained. Each bit is $X$-computable. Finite changes preserve density.
There is no claim of uniformly choosing a separator from arbitrary
endpoint descriptions.
\end{proof}

\begin{lemma}[A value beyond either preceding side]\label{lem:diagonal}
For every $k\geq1$ some $x\in\mathsf D_{k+1}\cap(0,1)$ lies in
neither $\mathsf S_k$ nor $\mathsf P_k$.
\end{lemma}
\begin{proof}
Enumerate all $\Sigma^0_k$ and $\Pi^0_k$ subsets of $\Q$, including
those that are not cuts. Starting with $[1/4,3/4]$, test its midpoint
against the next set with $0^{(k)}$. If it belongs, keep the left third;
otherwise keep the right third. If the tested set is the strict lower
cut of a value $v$, that third excludes $v$. The nested intervals have
length $3^{-n}/2$ and $0^{(k)}$-computable endpoints. Their intersection
is therefore in $\mathsf D_{k+1}$ and avoids every listed cut.
\end{proof}

Two qualifications will recur. Independent limits do not add their oracle
levels: $0'$ computes both limits of two computable sequences and their
defined computable combination. Additional levels arise from dependent
limits with the required uniform names. Also, a syntactic ceiling
does not prove a value has exactly that least level or that a physical
model attains it.
